\chapter{Discussion}
\label{ch:discussion}

\section{Physiological Heterogeneity, Not a Data-Volume Problem}

The bimodal shape of the LOPO result (Section~\ref{sec:results-lopo}, Figure~\ref{fig:lopo-bimodal}) is, in my view, the most theoretically interesting finding this project produced, and it is worth dwelling on rather than treating as a curiosity alongside the headline mean. If cross-patient generalisation failure were primarily a data-volume problem -- simply a matter of not having seen enough patients to interpolate smoothly between them -- one would expect sensitivity to degrade gradually as a function of some notion of "distance" from the training pool, producing a roughly unimodal distribution with a long tail. That is not what was observed. Sensitivity clusters sharply at the extremes, with only two of eight collapse-passing patients landing meaningfully between 0 and 1. This is more consistent with a picture in which the shared representation either does or does not happen to contain something structurally close enough to a given patient's own ictal morphology -- and when it does not, no amount of averaging across other patients rescues the result, because the problem is closer to a categorical mismatch than a continuous shortfall.

This account is corroborated, not merely asserted: chb03 and chb13's difficulty is independently flagged by \citet{shoeb2009}, on an entirely different detector, working with an entirely different methodology, some fifteen years before this project's own LOPO sweep identified the same two patients as structurally hard. Two detectors built with different tools, at different times, converging on the same patients being genuinely difficult is a stronger form of evidence than either result alone.

One caveat is worth stating honestly rather than left implicit. This project's preprocessing code contains a diagnostic note, and a corresponding but never-applied z-score 
normalisation step, describing a substantial amplitude discrepancy for chb03 between sessions. The project's actual amplitude-correction convention -- min-max scaling, tested directly 
with the amplitude-calibration experiment discussed above -- found no effect for chb03, which is the evidence the "structural, not amplitude-driven" claim rests on. 
Although to note, the specific alternative remedy noted in that preprocessing code (z-score normalisation with outlier clipping) was never itself tested end-to-end against chb03's result. The 
structural-difficulty account above is the settled explanation given the evidence actually collected, but this specific untested alternative is disclosed here rather 
than silently set aside.

\realfig{eeg_seizure_comparison_chb10_vs_chb13.png}{0.85}{Seizure-event waveform morphology compared between chb10 (a patient this pipeline handles well across every method tried) and chb13 (the project's consistently hardest patient), giving visual grounding to the repeated claim that chb13 is structurally, not merely statistically, difficult.}{fig:chb10-vs-chb13-waveform}

\section{Personalisation as the Architecturally Correct Response}

Taken together, the pool6-minus-X-versus-pool7 result (Section~\ref{sec:results-owndata}) and the C2 fine-tuning result (Section~\ref{sec:results-c2}) support the same underlying 
architectural conclusion with two different mechanisms: if inter-patient heterogeneity is genuinely closer to categorical than continuous, then a model's ability to access a given 
patient's own data -- whether via pool composition or via explicit fine-tuning -- is not a nice-to-have refinement but the architecturally appropriate response to the problem as it 
actually presents. The C2 chb13 case makes this concretely visible: fine-tuning did not simply raise a number, it corrected a specific failure \emph{mode} (over-triggering), which is a 
stronger legible claim than an isolated sensitivity improvement would have been.

It is worth returning explicitly, with concrete figures now in hand, to the triage-device framing set out in the Introduction (Section~\ref{sec:triage-framing}). This dissertation's 
unfiltered LOPO sensitivity (0.586) sits meaningfully below \citet{ali2024}'s unconstrained-compute comparator (72--75\%).  This gap is not explained away by the 
personalisation results above and it remains a real, disclosed shortfall in the cross-patient, zero-personalisation case. What the personalisation results \emph{do} establish is that this project's proposed clinical role is not the cold, shared-model number at all: a deployed device is not expected to operate as a frozen, one-size-fits-all classifier, but as a system that incorporates a given patient's own data over time, exactly the regime in which this project's strongest results (Sections~\ref{sec:results-owndata} and~\ref{sec:results-c2}) were obtained. Read this way, the triage framing and the personalisation results are mutually reinforcing rather than separate findings: a low-power, sensitivity-first screening tier is precisely the kind of system for which continued, low-cost personalisation on-device is both feasible and appropriate, in a way that matching an unconstrained-compute research system cold is not.

\section{Pool-Size and Own-Data Effects Are Genuinely Different Mechanisms}

A point worth stating explicitly, given how easily the two could be conflated in a less careful account: this project has now demonstrated \emph{two} distinct, genuine details that 
both involve pool7, and they answer different questions. The own-data-inclusion effect (Section~\ref{sec:results-owndata}) shows that a specific patient's own data, present in an 
otherwise near-constant-sized pool, measurably helps or clarifies an artefact for that specific patient in four of seven cases. The pool-size effect (Section~\ref{sec:results-c2}) 
shows, on a separate eight-patient set, that a larger pool alone -- independent of whether the evaluated patient's own data is in it -- substantially reduces the rate of 
collapse-artefact failures (37.5\% to 87.5\%), at a real cost in false-positive rate. Both are genuine, positive findings. Neither supersedes the other, and presenting them as a 
single, undifferentiated "more data helps" conclusion would understate the more precise and more useful thing actually learned: that at least two separable mechanisms 
are both at work, and a future system might benefit from exploiting them independently -- for instance, growing a shared background pool for general robustness while still prioritising 
acquisition of a specific patient's own calibration data for that patient's deployed device.

\section{Methodology as a Finding: The Caught-Bug Narrative}

In this project I reported several methodological errors caught and corrected during its own development, and this section treats that history as a connected narrative, 
rather than a list of incidents or embarrassments to be minimised. The chronological-slice-versus-LOPO-full conflation (Section~\ref{sec:eval-protocol}) silently produced a wrong LOPO 
re-run before being caught. An unused array construction in the pool-dataset-building script caused an out-of-memory failure, the second confirmed instance of an identical pattern 
already fixed once elsewhere in the same codebase. A shared default output path caused three of this project's C1 checkpoints to be silently overwritten mid-session, contaminating an 
earlier set of reported AUPRC figures that were subsequently retracted and rebuilt from verified weights. A figure-generation script was found to be reading from an entirely wrong 
file-naming pattern, silently mixing a hardcoded fallback value with live but methodologically mismatched data. The chb10 conformal-prediction reproducibility gap 
traced an unreproducible earlier figure to a script version lost to an uncommitted intermediate state. A sign-flip bug (Section~\ref{sec:results-hw}) 
appeared independently in both the simulator evaluation script and the separate hardware inference script.

The unifying pattern across several of these -- a recurring class of bug involving a shared, or wrongly assumed, default file path -- is now a recognised, named risk in this project's 
own codebase, not a collection of unrelated incidents. I state this pattern explicitly because it is, in my view, evidence \emph{for}  
how every one of these errors was caught and disclosed through my own verification practice and something that should not be omitted as a silent failure or embarrassment.

\section{The AUPRC-over-ROC-AUC Justification, Revisited}

The base-rate-relative AUPRC correction described in Section~\ref{sec:eval-protocol} is worth returning to here as a worked illustration of a more general methodological point, rather than a narrow technical footnote. Raw AUPRC, and ROC-AUC in isolation, can both present a misleadingly favourable or unfavourable picture on clinically imbalanced data unless corrected for each patient's own base rate. This project's own pre-registered expectation -- that chb15 would prove genuinely inert and chb13/chb16 would prove cases of simple misplaced thresholding -- was close to backwards once this correction was applied. A reader encountering only the raw AUPRC or ROC-AUC figures, without this correction, would likely have drawn the wrong conclusion about which patients were genuinely hard to discriminate versus merely poorly thresholded.

\section{Comparison to Ali et al. (2024) and Shoeb (2009)}

Both comparator systems introduced in Chapter~\ref{ch:background} are revisited here, with the full picture given in Table~\ref{tab:comparator-full} rather than a single headline gap figure, so the comparison can be judged in the round. Ali et al.'s figure is aligned in this table against this dissertation's raw, all-folds LOPO sensitivity rather than the collapse-corrected figure, since their own published methodology and per-subject results confirm they apply no equivalent artefact filter (Chapter~\ref{ch:background}) -- 
the collapse-diagnostic-filtered 0.470 is presented as a secondary analysis, not as this dissertation's headline figure, and is not the number set directly against theirs.
\begin{table}[H]
\centering
\small
\caption{Full comparison against both comparator systems, by evaluation regime. LOPO figures are compared against Ali et al.; per-patient figures are compared against Shoeb, since the two comparator systems answer different questions (Chapter~\ref{ch:background}).}
\label{tab:comparator-full}
\begin{tabular}{p{4cm}p{3cm}p{3cm}p{3cm}}
\toprule
Regime & This dissertation & Ali et al. (2024) & Shoeb (2009) \\
\midrule
LOPO event sensitivity (primary, unfiltered) & \textbf{0.59} ($\pm$0.45) & \textbf{72--75\%} & -- \\
\midrule
LOPO event sensitivity (secondary analysis, Section~\ref{sec:collapse-secondary}) & 0.47 ($\pm$0.44) & Not directly comparable & -- \\
\midrule
Per-patient (C2) mean sensitivity & \textbf{0.76} & -- & -- \\
\midrule
Per-patient detection latency & \textbf{as low as 1.0s} (default threshold) & -- & \textbf{4.6s mean} \\
\midrule
Hardware power/energy reported & \textbf{Yes} (Section~\ref{sec:results-hw}) & No & No \\
\midrule
Quantisation/hardware constraint & \textbf{Yes} (w4a4, AKD1000 v1) & No & No \\
\midrule
Collapse-diagnostic-equivalent screening & \textbf{Yes} & Not identified & Not identified \\
\bottomrule
\end{tabular}
\end{table}

Against \citet{ali2024}'s genuine LOPO comparator, this dissertation's primary, unfiltered sensitivity (0.586) falls short by approximately 0.15 against their midpoint (73.5\%) -- a gap, obtained on a fair, like-for-like basis, since neither system applies any artefact-filtering step to the figure being compared.
Against \citet{shoeb2009}'s per-patient comparator, this dissertation's fine-tuned per-patient latency (as low as 1.0 second at default threshold) matches or improves on Shoeb's reported mean of 4.6 seconds. Hardware is the one dimension along which this project is competitive without any comparability caveat at all: neither comparator system reports any hardware power or energy figure, since neither was designed under a comparable power constraint. This is stated plainly as this project's clearest genuine differentiator.

\section{Limitations}

Several honest limitations are stated here explicitly, consistent with this dissertation's standing practice throughout. chb16 fails this project's own false-positive-rate target under every calibration method tested (Section~\ref{sec:results-c2}), a genuine limitation for that specific patient rather than a resolved case in the sense of being fixed -- but the question of \emph{which} threshold to recommend for it has been resolved: the conformal-calibrated point ($\alpha=0.01$) is adopted as chb16's single recommended operating threshold, for consistency with every other C2 patient, superseding the two earlier, less rigorous candidates considered at prior stages of this project.

The hardware power measurement (Section~\ref{sec:results-hw}) carries an instrumentation limitation rather than an unresolved open decision: the system-level bench measurement cannot isolate the AKD1000 chip's own draw from Raspberry Pi host-CPU orchestration overhead, and true chip-isolated measurement via a shunt resistor was assessed and deliberately not attempted, given the risk of physical damage to the hardware this close to the project's deadline. Finally, several candidate techniques (seed ensembling, a sibling-script memory-pattern audit, Candidate G's ROC-AUC computation) remain logged as legitimate future directions rather than completed work, and are carried forward explicitly into Chapter~\ref{ch:conclusion} rather than silently dropped.

\section{Future Work: Formalising a Supplementary Diagnostic}
\label{sec:collapse-future-work}

One further limitation is stated here as its own section rather than folded into the list above, since it concerns my own methodology rather than a specific result. Throughout this project, I additionally screened results against a self-designed diagnostic intended to flag folds where an apparently high sensitivity is an artefact of near-permanent triggering rather than genuine discrimination. Since a model that fires on almost every window it sees will register a high event sensitivity by simple virtue of never staying silent through a genuine seizure. This diagnostic checks three explicit, numerical conditions; a fold or patient result is flagged if \emph{any} of the following trigger:

\begin{enumerate}
    \item \textbf{Window specificity below 0.9.} A model firing on more than 10\% of genuinely non-seizure windows is, in my judgement, already behaving closer to an alarm that is permanently or near-permanently active than to a discriminating detector, regardless of how high its reported sensitivity looks.
    \item \textbf{A single contiguous predicted block covering more than 90\% of the recording.} This catches a model that fires once and simply never stops, which technically overlaps every genuine seizure event (hence a sensitivity of 1.0) without ever having discriminated anything.
    \item \textbf{Positive-window fraction below 10\% of the true seizure-window rate.} The mirror-image failure -- a model that has effectively gone silent, firing so rarely that its few detections are more plausibly noise than genuine discrimination.
\end{enumerate}

\realfig{collapse_diagnostic_illustration.png}{0.85}{Three panels illustrating the supplementary diagnostic's three conditions: (a) a normal, discriminating prediction trace; (b) an over-firing collapse (a single near-permanent predicted block); (c) an under-firing collapse (almost no detections).}{fig:collapse-illustration}

\begin{verbatim}
def collapse_diagnostic(y_pred, window_specificity, y_true=None,
                         step_s=1.0, gap_tolerance_s=60.0,
                         spec_floor=0.9, block_frac_ceiling=0.9,
                         underfire_ratio_floor=0.1):
    # Condition (1): window specificity below the floor
    if window_specificity < spec_floor:
        reasons.append(f"window specificity {window_specificity} < {spec_floor}")
    # Condition (2): a single block covers too much of the recording
    if largest_block_frac > block_frac_ceiling:
        reasons.append("largest predicted block too large")
    # Condition (3): under-firing relative to the true seizure rate
    if pos_frac < underfire_ratio_floor * true_frac:
        reasons.append("positive-window fraction too low -- likely under-firing")
    return {"pass": len(reasons) == 0, "reasons": reasons, ...}
\end{verbatim}

I have deliberately not used this diagnostic to filter this dissertation's primary reported figures (Chapter~\ref{ch:results}), and it is presented here, in Limitations and Future Work, rather than in Chapter~\ref{ch:methodology} as though it were established methodology, for a specific reason: it is a heuristic I designed myself during this project, its three threshold values (0.9, 0.9, 0.1) reflect my own judgement rather than an externally validated standard, and neither \citet{ali2024} nor \citet{shoeb2009} implement anything directly equivalent against which it could be cross-checked. 
As such the future work for benchmarking this heuristic's three thresholds include being against a labelled dataset of known over-triggering and under-firing failures, comparing it to any equivalent screening steps used in the wider seizure-detection literature that this project has not identified, and determining whether its thresholds generalise beyond the specific architecture and dataset used in this project. Until that validation work is done, the secondary analysis in Section~\ref{sec:collapse-secondary} - including the bimodality finding - should be read as suggestive rather than as confirmed alongside this dissertation's primary results.